% Generated by build/sync_qft_soln.py from committed sources only.
% Source repository: https://github.com/Luca-Yucheng-Jin/QFT-soln
% Source commit: 07d7fe95fd5cf5b26d38a16ee535a04a925277b2
\section{Saddle-Point Approximation and Functional Determinants (PSI QFT II, Tutorial 3)}
\markboth{\thesection\quad SADDLE-POINT METHODS (PSI QFT II, TUTORIAL 3)}{}

\noindent\textit{Question prompts paraphrased from the official
\href{https://psi-online.perimeterinstitute.ca/courses/take/quantum-field-theory-ii-student/pdfs/20776373-module-3-tutorial-2}{Perimeter PSI Online tutorial sheet};
the equations and notation follow the source.}

\subsection{Saddle-Point Approximation}

Approximation methods for ordinary integrals are among the reasons that path
integrals are useful. For example, consider
\begin{equation}
    I(z)=\int dx\,e^{if(x,z)}.
\end{equation}
The corresponding ideas extend to functional integrals such as
\begin{equation}
    Z[J]=\int\mathcal D\phi\,e^{iS[\phi,J]}.
\end{equation}
The saddle-point method is one especially useful example.

Restore $\hbar$ in the Euclidean partition function:
\begin{equation}
    Z[J]
    =\int\mathcal D\phi\,
    e^{-\frac{1}{\hbar}S[\phi,J]}.
\end{equation}
This makes it possible to organise a semiclassical expansion by treating $\hbar$
as small. In Euclidean signature, configurations $\phi_c$ minimising the action
dominate because configurations with larger action are exponentially suppressed.
For a Minkowski integral, all stationary configurations of the action must be
considered rather than only its minima.

The following questions first develop this approximation for ordinary
$n$-dimensional integrals.

\begin{enumerate}[label=\alph*)]
    \item Let $S(x^\mu;z)$ be a real function of $n$ spacetime variables
    $\mu=1,\ldots,n$ and a coupling $z$, and define
    \begin{equation}
        I(z)
        =\int_{-\infty}^{\infty}d^n x\,
        e^{-\frac{1}{\hbar}S(x;z)}.
    \end{equation}
    Expand $S(x;z)$ around its extrema $x_c$ and show that the sum over its
    minima gives
    \begin{equation}
        I(z)
        \approx
        \sum_{\operatorname{Min}[x_c]}
        e^{-\frac{1}{\hbar}S(x_c;z)}
        \left[
            \frac{(2\pi\hbar)^n}{\det S''(x_c)}
        \right]^{1/2}
        +\mathcal O(\hbar^a),
    \end{equation}
    where
    \begin{equation}
        \left[S''(x_c)\right]_{\mu\nu}
        =\left.
        \frac{\partial^2S}{\partial x^\mu\,\partial x^\nu}
        \right|_{x=x_c}.
    \end{equation}
    \begin{solution}
        Expanding $S(x;z)$ around its minima, we find
        \begin{equation}
            S(x;z) = S(x_c;z) + \frac{1}{2}(x-x_c)^\mu (x-x_c)^\nu [S''(x_c;z)]_{\mu\nu} + \cdots.
        \end{equation}
        In the semi-classical regime, the field configurations around the minima dominate, giving
        \begin{equation}
        \begin{aligned}
            I(z) &\approx \sum_{\operatorname{Min}[x_c]}\int d^nx e^{-\frac{1}{\hbar}(S(x_c;z) + \frac{1}{2}(x-x_c)^\mu (x-x_c)^\nu [S''(x_c;z)]_{\mu\nu})}\\
            &=\sum_{\operatorname{Min}[x_c]} e^{-\frac{1}{\hbar}S(x_c;z)}\left[
            \frac{(2\pi\hbar)^n}{\det S''(x_c)}
        \right]^{1/2},
        \end{aligned}
        \end{equation}
        where we have used the results from conventional Gaussian integration in the last equality.
    \end{solution}

    \item For a generic function $S(x;z)$, determine the exponent $a$.
    \begin{solution}
        The important fluctuations should be $x-x_c \sim \sqrt{\hbar}$. The next contribution comes from $S''''$ ($S'''$ is odd, and hence vanishes). Therefore, $a = n/2 +2 - 1 = n/2 +1$.
    \end{solution}

    \item State the condition under which the approximation becomes exact.
    \begin{solution}
        The approximation is exact if all higher order derivative vanishes.
    \end{solution}

    \item Explain how the saddle-point expression should change if the starting
    point is the Minkowski partition function
    \begin{equation}
        Z=\int e^{iS/\hbar}.
    \end{equation}
    \begin{solution}
        In real time, $-S \to iS$. Besides that, we also have to add the $i\epsilon$ prescription to ensure the path integral converges. Therefore, we have
        \begin{equation}
            I(z)\approx
        \sum_{\operatorname{Min}[x_c]}
        e^{\frac{i}{\hbar}S(x_c;z)}
        \left[
            \frac{(2\pi i\hbar)^n}{\det (S''(x_c) + i\epsilon)}
        \right]^{1/2}
        +\mathcal O(\hbar^a).
        \end{equation}
    \end{solution}
\end{enumerate}

\subsection{Functional Determinants}

Functional determinants occur throughout the path-integral treatment of QFT,
including the semiclassical result above. Their formal definition as a product
over operator eigenvalues is generally divergent and requires regularisation.
For this exercise, set those analytic subtleties aside and manipulate the
operators by analogy with matrices.

Define
\begin{equation}
    F[f]
    =\log\!\left(
        \det\!\left[-\frac{d^2}{dy^2}+f(y)\right]
    \right).
\end{equation}
Thus $F$ is the logarithm of the determinant of the operator
\begin{equation}
    \mathcal A[f(y)]
    =-\frac{d^2}{dy^2}+f(y).
\end{equation}
If $M(y,z)$ is the integral kernel of an operator $\mathcal M$, then
\begin{equation}
    \mathcal M g(y)
    =\int M(y,z)g(z)\,dz
\end{equation}
for every suitable function $g$.

The matrix-product analogy used here is
\begin{equation}
    \sum_y M_{xy}M_{yz}
    \quad\Longleftrightarrow\quad
    \int dy\,M(x,y)M(y,z).
\end{equation}

\begin{enumerate}[label=\alph*)]
    \item Show that the integral kernel of $\mathcal A$ is
    \begin{equation}
        A(y,z)
        =\left[-\frac{d^2}{dy^2}+f(y)\right]\delta(y-z).
    \end{equation}
    \begin{solution}
        \begin{equation}
            \begin{aligned}
                \mathcal{A}[f(y)]g(y) = -\frac{d^2g}{dy^2} + fg = \int dz \, A(y,z) g(z).
            \end{aligned}
        \end{equation}
    \end{solution}

    Suppose that $O(x,y)$ is the integral kernel of
    \begin{equation}
        \mathcal O
        =\left(-\frac{d^2}{dy^2}+f(y)\right)^{-1}.
    \end{equation}

    \item The operators satisfy $\mathcal A\mathcal O=1$. Rewrite this identity
    entirely in terms of their integral kernels.
    \begin{solution}
        \begin{equation}
            \begin{aligned}
                \int dz A(x,z)O(z,y) = \delta(x-y).
            \end{aligned}
        \end{equation}
    \end{solution}

    \item Using functional differentiation, prove
    \begin{equation}
        \frac{\delta F[f]}{\delta f(x)}=O(x,x).
    \end{equation}
    \textit{Hint:} Recall the matrix identity
    \begin{equation}
        \det M=\exp(\operatorname{tr}\log M).
    \end{equation}
    \begin{solution}
        \begin{equation}
            \begin{aligned}
                \frac{\delta F[f]}{\delta f(x)} &= \frac{\delta}{\delta f(x)} \operatorname{tr} \log \mathcal{A} = \frac{\delta}{\delta f(x)}\operatorname{tr} \mathcal{O} \delta\mathcal{A}= \frac{\delta}{\delta f(z)} \int dydz O(x,y)\delta f(y) \delta(y-z) =O(x,x).
            \end{aligned}
        \end{equation}
    \end{solution}

    \item Next show
    \begin{equation}
        \frac{\delta^2F[f]}{\delta f(x)\delta f(y)}
        =-O(x,y)O(y,x).
    \end{equation}
    \textit{Hint:} To differentiate an inverse, begin with
    \begin{equation}
        M^{-1}M=1.
    \end{equation}
    \begin{solution}
        \begin{equation}
            \frac{\delta^2F[f]}{\delta f(x)\delta f(y)} = \frac{\delta}{\delta f(y)}\operatorname{tr}\mathcal{A}^{-1} = -\operatorname{tr}\mathcal{A}^{-1}\frac{\delta\mathcal{A}}{\delta f(y)}\mathcal{A}^{-1} = - O(x,y) O(y,x).
        \end{equation}
    \end{solution}

    \item Finally derive
    \begin{align}
        \frac{\delta^3F[f]}
        {\delta f(x)\delta f(y)\delta f(z)}
        &=O(x,y)O(y,z)O(z,x)
        \notag\\
        &\quad+O(x,z)O(z,y)O(y,x).
    \end{align}
    \begin{solution}
        Follows trivially from above.
    \end{solution}
\end{enumerate}

\subsection{Optional: The Gel'fand--Yaglom Formalism}

The Gel'fand--Yaglom method can determine functional determinants without
explicitly computing every eigenvalue. In favourable cases it is exact, and it
can also be implemented numerically. Begin with the Riemann zeta function
\begin{equation}
    \zeta(s)
    \equiv
    \sum_{k=1}^{\infty}k^{-s}.
\end{equation}

The contour required below is fixed by the argument principle.  It is positively
oriented and surrounds the positive zeros $\lambda_n$ of $F(\lambda)$, while
avoiding the branch cut of $\lambda^{-s}$ on the negative real axis.

\begin{figure}[H]
    \centering
    \begin{tikzpicture}[
        >=Stealth,
        scale=0.95,
        every node/.style={font=\small},
        gyAxis/.style={->,draw=black,line width=0.45pt},
        gyContour/.style={
            draw=black,
            line width=0.9pt,
            line cap=round,
            line join=round,
            postaction={decorate},
            decoration={markings,
                mark=at position 0.18 with
                    {\arrow{Stealth[length=2.1mm]}},
                mark=at position 0.67 with
                    {\arrow{Stealth[length=2.1mm]}}
            }
        },
        gyCut/.style={draw=black,line width=0.55pt,
                     line cap=round,line join=round}]
        \draw[gyAxis] (-4.2,0) -- (4.8,0)
            node[right] {$\operatorname{Re}\lambda$};
        \draw[gyAxis] (0,-1.85) -- (0,1.85)
            node[above] {$\operatorname{Im}\lambda$};

        \foreach \x in {-4.00,-3.75,...,-0.25}
            \draw[gyCut] (\x,0) -- ++(0.125,0.055) -- ++(0.125,-0.055);
        \node[font=\scriptsize,fill=white,inner sep=1pt] at (-2.05,0.42)
            {branch cut};
        \draw[-{Stealth[length=1.4mm]},line width=0.35pt]
            (-2.05,0.30) -- (-1.92,0.08);
        \node[below left] at (0,0) {$0$};

        \foreach \x/\lab in {
            0.85/{\lambda_1},
            1.65/{\lambda_2},
            2.45/{\lambda_3}}
        {
            \fill (\x,0) circle (1.6pt);
            \node[below=3pt] at (\x,0) {$\lab$};
        }
        \node at (3.20,0) {$\cdots$};

        \draw[gyContour]
            (0.35,-1.12)
            -- (3.65,-1.12)
            .. controls (4.25,-1.12) and (4.25,1.12) ..
            (3.65,1.12)
            -- (0.35,1.12)
            .. controls (-0.02,0.78) and (-0.02,-0.78) .. cycle;
        \node at (3.92,1.28) {$\gamma$};
    \end{tikzpicture}
    \caption{Schematic reconstruction of contour $\gamma$
    in the $\lambda$-plane.}
\end{figure}

Let $A$ have eigenvalues $\lambda_n$ and eigenfunctions $\phi_n$:
\begin{equation}
    A\phi_n=\lambda_n\phi_n.
\end{equation}
Associate to it the zeta function
\begin{align}
    \zeta_A(s)
    &\equiv
    \operatorname{Tr}\!\left[\frac{1}{A^s}\right]
    \notag\\
    &=\sum_{k=1}^{\infty}\lambda_k^{-s}.
\end{align}

\begin{enumerate}[label=\alph*)]
    \item Show that zeta regularisation gives
    \begin{equation}
        \det A=\exp\!\left[-\zeta_A'(0)\right].
    \end{equation}
    Because $\zeta_A$ and its derivative need not converge near $s=0$,
    analytic continuation may be required before this expression is meaningful.
    \begin{solution}
        \begin{equation}
            \det A =\prod_{k=1}^\infty \lambda_k,
        \end{equation}
        while
        \begin{equation}
            \zeta'_A = \sum_{k=1}^\infty-\log \lambda_k \lambda^{-s}.
        \end{equation}
        Therefore,
        \begin{equation}
            \det A=\exp\!\left[-\zeta_A'(0)\right].
        \end{equation}
    \end{solution}

    \item Take
    \begin{equation}
        A=-\frac{d^2}{d\tau^2}
    \end{equation}
    on $0\leq\tau\leq1$, with Dirichlet conditions
    \begin{equation}
        \phi_n(0)=\phi_n(1)=0.
    \end{equation}
    Find the eigenvalues $\lambda_n$, then evaluate $\det A$ using
    \begin{equation}
        \zeta(0)=-\frac{1}{2},
        \qquad
        \zeta'(0)=-\frac{1}{2}\log(2\pi).
    \end{equation}
    \begin{solution}
        It can be easily found that $\lambda_n = n^2\pi^2$, $n \in \mathbb{Z}^{+}$. Then,
        \begin{equation}
            \zeta_A (s) = \sum_{k=1}^\infty(n\pi)^{-2s},
        \end{equation}
        and
        \begin{equation}
            \zeta_A'(0) = \sum_{k=1}^\infty-2\log k - 2\log\pi = 2\zeta'(0) -2\log \pi\zeta(0) = -\log2.
        \end{equation}
        Hence, we have
        \begin{equation}
            \det A = 2.
        \end{equation}
    \end{solution}

    In cases where the eigenvalues cannot be obtained explicitly, suppose one
    can construct $F(\lambda)$ whose zeros are precisely the eigenvalues of $A$.

    \item Assume that the $\lambda_n>0$ are discrete and nondegenerate. Show that
    \begin{equation}
        \zeta_A(s)
        =\frac{1}{2\pi i}
        \int_\gamma d\lambda\,
        \lambda^{-s}
        \frac{d\ln F(\lambda)}{d\lambda},
    \end{equation}
    where $\gamma$ is the contour specified in Figure 1 and the branch cut lies
    on the negative real axis.
    \begin{solution}
        The RHS is given by
        \begin{equation}
            \begin{aligned}
                &\frac{1}{2\pi i}
        \int_\gamma d\lambda\,
        \lambda^{-s}
        \frac{d\ln F(\lambda)}{d\lambda}\\
        &= \sum_{n=1}^\infty
        \lim_{\lambda\to\lambda_n} \frac{\lambda-\lambda_n}{\lambda^s}\frac{1}{F(\lambda)}\frac{d F(\lambda)}{d\lambda}\\
        &=\sum_{n=1}^\infty
        \lim_{\lambda\to\lambda_n}\left(F'(\lambda_n) + F''(\lambda_n)(\lambda-\lambda_n) + \cdots \right)\left(F'(\lambda_n)(\lambda-\lambda_n) + \frac12F''(\lambda_n)(\lambda-\lambda_n)^2 + \cdots \right)^{-1}\\
        &=\sum_{n=1}^\infty
        \lim_{\lambda\to\lambda_n}\frac{1}{\lambda^s}\left(1+ \frac{F''(\lambda_n)}{F'(\lambda_n)}(\lambda-\lambda_n) +\cdots\right)\left(1- \frac12\frac{F''(\lambda_n)}{F'(\lambda_n)}(\lambda-\lambda_n) +\cdots\right)\\
        &=\sum_{n=1}^\infty \lambda_n^{-s}\\
        &= \zeta_A(s).
            \end{aligned}
        \end{equation}
    \end{solution}

    \item Deform $\gamma$ into the contour $\gamma_-$, which hugs the negative
    real axis as indicated in Figure 2, and derive
    \begin{equation}
        \zeta_A(s)
        =\frac{\sin(\pi s)}{\pi}
        \int_0^{-\infty}d\lambda\,
        \lambda^{-s}
        \frac{d\ln F(\lambda)}{d\lambda}.
    \end{equation}
    \begin{solution}
        \begin{equation}
        \begin{aligned}
                        \zeta_A(s) &= \frac{1}{2\pi i}
        \int_{\gamma_-} d\lambda\,
        \lambda^{-s}
        \frac{d\ln F(\lambda)}{d\lambda}\\ &= \lim_{\epsilon \to 0^+}\frac{1}{2\pi i}
        \int_{-\infty}^0 d\lambda\,
        (\lambda+i\epsilon)^{-s}
        \frac{d\ln F(\lambda+i\epsilon)}{d\lambda} -  \frac{1}{2\pi i}
        \int_{-\infty}^0 d\lambda\,
        (\lambda-i\epsilon)^{-s}
        \frac{d\ln F(\lambda-i\epsilon)}{d\lambda}\\
        &=\frac{1}{2\pi i}
        \int_{-\infty}*0 d\lambda\,
        (e^{-s\log \lambda-is\pi}-e^{-s\log \lambda+ is\pi})
        \frac{d\ln F(\lambda)}{d\lambda}\\
        &=-\frac{\sin \pi s}{\pi}\int_{-\infty}^0 d\lambda\,
        \lambda^{-s}
        \frac{d\ln F(\lambda)}{d\lambda}.
        \end{aligned}
        \end{equation}
    \end{solution}


    \begin{figure}[H]
        \centering
        \begin{tikzpicture}[
            >=Stealth,
            scale=0.82,
            every node/.style={font=\small},
            gyAxis/.style={->,draw=black,line width=0.45pt},
            gyOriginal/.style={
                draw=black!48,
                densely dashed,
                line width=0.65pt,
                line cap=round,
                line join=round,
                postaction={decorate},
                decoration={markings,
                    mark=at position 0.18 with
                        {\arrow{Stealth[length=1.8mm]}},
                    mark=at position 0.67 with
                        {\arrow{Stealth[length=1.8mm]}}
                }
            },
            gyBank/.style={
                draw=black,
                line width=0.9pt,
                line cap=round,
                postaction={decorate},
                decoration={markings,
                    mark=at position 0.56 with
                        {\arrow{Stealth[length=2.1mm]}}
                }
            },
            gyOuter/.style={
                draw=black,
                line width=0.9pt,
                line cap=round,
                line join=round,
                postaction={decorate},
                decoration={markings,
                    mark=at position 0.50 with
                        {\arrow{Stealth[length=2.1mm]}}
                }
            },
            gyCurve/.style={draw=black,line width=0.9pt,
                           line cap=round,line join=round},
            gyCut/.style={draw=black,line width=0.55pt,
                         line cap=round,line join=round}]
            \draw[gyAxis] (-5.15,0) -- (5.25,0)
                node[right] {$\operatorname{Re}\lambda$};
            \draw[gyAxis] (0,-2.45) -- (0,2.45)
                node[above] {$\operatorname{Im}\lambda$};

            \foreach \x in {-4.50,-4.25,...,-0.25}
                \draw[gyCut] (\x,0) -- ++(0.125,0.055)
                    -- ++(0.125,-0.055);
            \node[font=\scriptsize,fill=white,inner sep=1pt]
                at (-3.25,0.33) {branch cut};
            \draw[-{Stealth[length=1.4mm]},line width=0.35pt]
                (-3.25,0.22) -- (-3.10,0.07);
            \node[below right] at (0,0) {$0$};

            \foreach \x/\lab in {
                0.85/{\lambda_1},
                1.65/{\lambda_2},
                2.45/{\lambda_3}}
            {
                \fill (\x,0) circle (1.6pt);
                \node[below=3pt] at (\x,0) {$\lab$};
            }
            \node at (3.18,0) {$\cdots$};

            \draw[gyOriginal]
                (0.35,-0.78)
                -- (3.55,-0.78)
                .. controls (4.05,-0.78) and (4.05,0.78) ..
                (3.55,0.78)
                -- (0.35,0.78)
                .. controls (-0.02,0.55) and (-0.02,-0.55) .. cycle;
            \node[black!55] at (3.82,0.96) {$\gamma$};

            \draw[gyBank] (-4.20,0.50) -- (-0.42,0.50);
            \draw[gyCurve]
                (-0.42,0.50)
                .. controls (0.18,0.48) and (0.18,-0.48) ..
                (-0.42,-0.50);
            \draw[gyBank] (-0.42,-0.50) -- (-4.20,-0.50);
            \draw[gyOuter]
                (-4.20,-0.50)
                .. controls (-4.85,-0.82) and (-4.65,-1.62) ..
                (-3.70,-1.93)
                .. controls (-1.10,-2.62) and (3.55,-2.38) ..
                (4.38,-0.65)
                .. controls (4.65,-0.22) and (4.65,0.22) ..
                (4.38,0.65)
                .. controls (3.55,2.38) and (-1.10,2.62) ..
                (-3.70,1.93)
                .. controls (-4.65,1.62) and (-4.85,0.82) ..
                (-4.20,0.50);

            \node[above] at (-3.72,0.50) {$\gamma_-$};
        \end{tikzpicture}
        \caption{Finite contour deformation.}
    \end{figure}

    \item Apply the preceding results to the one-dimensional Schr\"odinger
    operator
    \begin{equation}
        A=-\frac{d^2}{dx^2}+V(x)
    \end{equation}
    on $0\leq x\leq1$, with Dirichlet boundary conditions:
    \begin{equation}
        \begin{aligned}
            A\phi_n&=\lambda_n\phi_n,\\
            \phi(0)&=0,\\
            \phi(1)&=0.
        \end{aligned}
    \end{equation}
    Assume that the $\lambda_n$ are positive, discrete, and nondegenerate.
    Instead of the Dirichlet eigenfunctions, solve the same eigenvalue equation
    with initial-value conditions
    \begin{equation}
        \begin{aligned}
            Au_\lambda&=\lambda u_\lambda,\\
            u_\lambda(0)&=0,\\
            u_\lambda'(0)&=1.
        \end{aligned}
    \end{equation}
    Choose
    \begin{equation}
        F(\lambda)=u_\lambda(1)
    \end{equation}
    and use the previous parts to establish the Gel'fand--Yaglom theorem
    \begin{equation}
        \frac{\det A}{\det A_{\mathrm{free}}}
        =\frac{u_0(1)}{u_{\mathrm{free},0}(1)},
    \end{equation}
    where
    \begin{equation}
        A_{\mathrm{free}}=-\frac{d^2}{dx^2}.
    \end{equation}
    \begin{solution}
        From the results above, we have
        \begin{equation}
            \zeta_A'(0) = \int^{-\infty}_0d\lambda \frac{d\log F(\lambda)}{d\lambda} = \log\frac{F(-\infty)}{F(0)} \implies \det A = \exp [-\zeta_A'(0)] = \frac{F(0)}{F(-\infty)}.
        \end{equation}
        Note that $F(-\infty) = u_{-\infty}(1)$ is formally infinite. Intuitively, one might guess
        \begin{equation}
            \lim_{\lambda\to -\infty}\frac{F(\lambda)}{F_\mathrm{free}(\lambda)} = 1,
        \end{equation}
        which follows
        \begin{equation}
            \frac{\det A}{\det A_{\mathrm{free}}}
        =\frac{u_0(1)}{u_{\mathrm{free},0}(1)}.
        \end{equation}
        The guess is true, as long as $V(x)$ is bounded, when $\lambda \to -\infty$, $V(x)$ becomes negligible.
    \end{solution}

    \item Use the Gel'fand--Yaglom theorem to compute $\det A$ for
    \begin{equation}
        A=-\frac{d^2}{dx^2}+\omega^2
    \end{equation}
    on the unit interval with Dirichlet boundary conditions. Verify the result
    independently by finding the eigenvalues and using
    \begin{equation}
        \sin(\pi z)
        =\pi z\prod_{n=1}^{\infty}
        \left(1-\frac{z^2}{n^2}\right).
    \end{equation}
    \begin{solution}
        \begin{equation}
            u_{\mathrm{free},0}(x) = x, \quad u_0(x) = \frac{\sinh \omega x}{\omega}, \quad \det A_\mathrm{free} = 2 \implies \det A =2\frac{\sinh \omega}{\omega}.
        \end{equation}
        Meanwhile,
        \begin{equation}
            \lambda_n =  \pi^2n^2 + \omega^2 \implies \det A = \prod_{n=1}^\infty \pi^2n^2\left(1 +\frac{\omega^2}{\pi^2 n^2}\right),
        \end{equation}
        which looks disastrous. However, remember, we can regulate it with
        \begin{equation}
            \det A_\mathrm{free} = \prod_{n=1}^\infty \pi^2n^2.
        \end{equation}
        Therefore, we find
        \begin{equation}
            \frac{\det A}{\det A_\mathrm{free}}  =\frac{\sinh \omega}{\omega},
        \end{equation}
        which recovers the Fel'danf-Yaglom theorem result by regulating $\det A_\mathrm{free}$ through the analytic continuation of the Reimann function.
    \end{solution}
\end{enumerate}
